\documentclass[11pt]{article}
\usepackage[margin=1in]{geometry}
\usepackage{times}
\usepackage[T1]{fontenc}
\usepackage[utf8]{inputenc}
\usepackage{enumitem}
\usepackage{xcolor}
\usepackage{hyperref}

\definecolor{reviewercolor}{RGB}{0,0,140}
\definecolor{responsecolor}{RGB}{0,100,0}

\newcommand{\reviewer}[1]{\textcolor{reviewercolor}{\textbf{Reviewer comment:} \textit{#1}}}
\newcommand{\response}[1]{\textcolor{responsecolor}{\textbf{Response:} #1}}
\newcommand{\location}[1]{\textbf{Location:} #1}

\title{Author Response --- Round 6}
\author{}
\date{}

\begin{document}
\maketitle

We thank the reviewers for their detailed and constructive feedback. We have revised the paper to address each concern. Below, we respond point-by-point, noting specific changes and their locations.

\section*{Overview of Changes}

\begin{itemize}[leftmargin=*]
\item New title scoped to observed evidence rather than general claims.
\item Introduction, Abstract, and Conclusion rewritten to multiplicity-first framing: per-example multiplicity ($\alpha = 0.14$--$0.30$, Jaccard $\leq 0.10$) now leads; convergence introduced as contrastive secondary finding.
\item Introduction findings reordered (First = multiplicity, Second = convergence); transition sentence explains why Results section retains convergence-first order (\S4.1 provides aggregate backdrop for \S4.2 per-example analysis).
\item All instances of ``general property'' replaced with empirically scoped language.
\item Per-example analysis expanded from two to four benchmarks (LAMBADA, HellaSwag, PIQA, ARC-Easy) plus BoolQ and WinoGrande as supplementary in Appendix; multiplicity spectrum $\alpha = 0.14$--$0.30$.
\item Bootstrap 95\% CIs on $\alpha$ added (Appendix Table~6 + \S4.2): all adjacent-benchmark CIs non-overlapping except PIQA--HellaSwag.
\item 340M LAMBADA per-example analysis added (Appendix \S F + \S4.3): $\alpha = 0.31$, $\bar{J} \approx 0.09$, confirming cross-scale stability of multiplicity.
\item Subsample sensitivity of $\alpha$ added (Appendix \S G + Method \S2.2): CI width at $n = 100$ spans 0.13--0.18, justifying COPA exclusion.
\item Cross-benchmark confidence margin analysis added (Appendix): rescue pattern replicates 9/9 on main benchmarks; disruption gradient quantified via coefficient of variation.
\item Table~2: $\rho_\sigma$ column renamed to $\rho_\alpha$ (Spearman between attention ratio and ambiguity $\alpha$); LAMBADA value corrected from $-1.0$ to $-0.6$.
\item \S4.4 expanded to detail the GDN head-count confound in the main text.
\item ARR Responsible NLP Research Checklist added.
\item Table~1 footnote explains the $p$-value rounding discrepancy.
\item Limitations section strengthened with single-checkpoint caveat.
\item \texttt{wang2025hybrid} bib entry corrected (title, authors).
\end{itemize}

\section*{Detailed Responses}

\subsection*{W-CITE (R2-N1): Incorrect citation for model zoo}

\reviewer{The citation for the model zoo (wang2025hybrid) has incorrect title and author list.}

\response{Fixed. The bib entry now reads ``A Systematic Analysis of Hybrid Linear Attention'' with authors Dustin Wang, Rui-Jie Zhu, Katharopoulos, and Gu.}

\location{\texttt{references.bib}, entry \texttt{wang2025hybrid}.}

\subsection*{W-CHECKLIST (R2-N2): Missing ARR Responsible NLP Research Checklist}

\reviewer{The submission lacks the mandatory ARR Responsible NLP Research Checklist.}

\response{Added. The checklist covers all sections (A1--A4, B1--B6, C1--C3, D, E) with specific justifications. Notably, C1 reports the computational budget (single NVIDIA RTX PRO 6000, $\sim$5 GPU-hours inference-only) and A3 explicitly notes the scoped claims.}

\location{\texttt{arr\_checklist.tex}, included after the Ethics Statement in \texttt{main.tex}.}

\subsection*{W-OVERCLAIM (R1-W1, R2-N3): Overclaiming on convergence}

\reviewer{The 79\% reduction is presented prominently without immediately noting that the adjusted p-value does not reach significance. The title claims ``Task-Specific Convergence'' which may overstate the strength of evidence.}

\response{Four changes address this:}
\begin{enumerate}[leftmargin=2em]
\item \textbf{Title}: Changed from ``Task-Specific Convergence and Per-Example Predictive Multiplicity in Hybrid Recurrent Architectures'' to ``Per-Example Predictive Multiplicity in Hybrid Recurrent Architectures.'' The convergence claim is removed from the title.
\item \textbf{Abstract}: Rewritten to lead with per-example multiplicity ($\alpha = 0.14$--$0.30$, Jaccard $\leq 0.10$); convergence introduced with ``by contrast'' as secondary finding.
\item \textbf{Introduction}: Opening hook rewritten from convergence to per-example multiplicity (Jaccard $\leq 0.10$ is the first number readers encounter); convergence demoted to background in sentence~2. Findings reordered: First = multiplicity spectrum, Second = convergence.
\item \textbf{Conclusion}: Likewise rewritten to multiplicity-first; convergence follows with ``by contrast.''
\end{enumerate}

\location{Title (line 1 of \texttt{main.tex}); Abstract (paragraph 1); Introduction \S1 (paragraphs 1 and 5); Conclusion \S5.}

\subsection*{W-GENERAL (R1-W2): ``General property'' overclaims from limited evidence}

\reviewer{Claiming per-example multiplicity is a ``general property of hybrid architectures'' is too strong given only three architectures and two benchmarks.}

\response{We have replaced all instances of ``general property'' with empirically scoped language and expanded the per-example analysis from two to four benchmarks (plus two supplementary):}
\begin{itemize}[leftmargin=2em]
\item Per-example analysis now covers four benchmarks spanning three task categories: LAMBADA (retrieval), HellaSwag (reasoning), PIQA (reasoning), ARC-Easy (mixed). BoolQ and WinoGrande are reported as supplementary in the Appendix due to structural limitations affecting Jaccard significance.
\item Multiplicity is characterized as a task-dependent spectrum: $\alpha$ ranges from 0.14 (PIQA) to 0.30 (LAMBADA).
\item PIQA and HellaSwag exhibit the lowest multiplicity ($\alpha = 0.14$--$0.15$), confirming that multiplicity magnitude is task-dependent rather than noise.
\item Within-benchmark $\alpha$ stability verified for all four benchmarks: range 0.004 (HellaSwag), 0.011 (PIQA), 0.025 (ARC-Easy), 0.022 (LAMBADA), all below 0.03 (Appendix, Table~11).
\item Conclusion ensemble claim: now qualified with ``if confirmed in additional architectures''
\end{itemize}

\location{Abstract, \S1, \S3.4, \S4.3 (Table~2), \S4.4, \S5, Limitations, Appendix.}

\subsection*{W-GDN (R1-W4): GDN head-count confound}

\reviewer{GDN's h3-1 configuration uses 16 attention heads rather than 32, creating a confound. The sensitivity analysis is only in the appendix.}

\response{\S4.4 (Robustness) now contains two paragraphs. The first retains the existing sensitivity test: removing GDN weakens $\rho$ from $-1.0$ to $-0.7$ ($p = .233$), while the RetNet--HGRN2 pair preserves the decoupling (all hybrid $\sigma \leq 0.008$, Jaccard saturates at 0.086--0.108). The second paragraph, newly added, explains the mechanism: GDN's h3-1 uses 16 attention heads vs.\ 32 in all other configurations, introducing an attention-capacity asymmetry. We state the conclusion explicitly: ``The aggregate convergence result is therefore partially dependent on GDN; the decoupling result is not.''}

\location{\S4.4, paragraph 2.}

\subsection*{R3-POWER (R1-W3): Low statistical power ($n = 5$)}

\reviewer{With $n = 5$ attention ratios, the minimum achievable p-value is $\sim$0.008, limiting the power of the convergence test.}

\response{We acknowledge this as a structural constraint of the model zoo design. It is discussed in three places: (1) the Method section describes the exact permutation test with 120 permutations; (2) \S4.1 reports the BCa bootstrap CI $[-1.0, -0.7]$ as a complementary characterization alongside the $p$-value; (3) the Limitations section (``Statistical power'' paragraph) explicitly states the minimum achievable $p$-value ceiling and its consequences. We agree that this limits the strength of the convergence finding, which is why we have removed convergence from the title and weakened all convergence-related claims (see W-OVERCLAIM above). The per-example multiplicity finding, which is the paper's central contribution, does not depend on the convergence $p$-value. Additionally, bootstrap 95\% CIs on $\alpha$ confirm the multiplicity spectrum ordering (all adjacent-benchmark CIs non-overlapping except PIQA--HellaSwag), and subsample sensitivity analysis shows $\alpha$ stabilizes at $n \geq 500$ (CI width $< 0.08$; Appendix~\S G).}

\location{\S3.5, \S4.1 (paragraph 1), Limitations (``Statistical power'' paragraph), Appendix~\S G.}

\subsection*{R3-BINARY (R1-W5): Binary correct/incorrect decomposition}

\reviewer{The helped/hurt decomposition uses binary correctness, which may lose information compared to continuous confidence scores.}

\response{The binary decomposition is a deliberate methodological choice aligned with the predictive multiplicity framework of Marx et al.\ (2020), which defines multiplicity over discrete prediction outcomes. Additionally, we report continuous-valued analyses that are consistent with the binary results: Cohen's $d$ for pure-model confidence margins ranges from 0.73 to 0.92 between helped and hurt examples on LAMBADA (Appendix, ``Confidence Margin Analysis''), showing that the binary decomposition captures a real separation in the continuous confidence space. We have extended this analysis to all four benchmarks:}
\begin{itemize}[leftmargin=2em]
\item \textbf{Rescue pattern (C2)}: Near-miss errors rescued by attention replicate across all three main benchmarks with confidence margin data (9/9 tests significant; $d = 0.13$--$0.17$ on HellaSwag under \texttt{acc\_norm}, $d = 0.51$--$0.60$ on ARC-Easy, $d = 0.73$--$0.77$ on LAMBADA). On BoolQ (binary classification), the rescue signal is weaker but directionally consistent (2/3, $d = 0.16$--$0.36$).
\item \textbf{Disruption pattern (C1)}: Effect sizes decrease with answer-space dimensionality: open-ended prediction ($d = 0.83$--$0.92$, LAMBADA) $>$ ARC-Easy ($d = 0.13$--$0.26$) $>$ HellaSwag ($d \approx 0$ under \texttt{acc\_norm}) $\approx$ binary ($d \approx 0$, BoolQ). This gradient broadly tracks the coefficient of variation (CV) of the log-likelihood confidence proxy: LAMBADA (CV $= 0.75$) $>$ ARC-Easy ($0.54$) $>$ HellaSwag ($0.49$) $>$ BoolQ ($0.22$). BoolQ's CV is $3.4\times$ lower than LAMBADA's, leaving insufficient relative spread to distinguish borderline-correct from robustly-correct examples. The asymmetry with C2 is expected: a correct answer in multiple-choice settings can be robust despite low absolute confidence if all alternatives score worse, whereas near-miss errors remain identifiable by relatively high target probability regardless of task format.
\end{itemize}
Critically, the multiplicity findings themselves---ambiguity $\alpha$ and pairwise Jaccard---do not depend on the confidence margin analysis. Multiplicity is present and statistically significant on all four benchmarks including BoolQ ($\alpha = 0.45$, 12/24 Jaccard pairs significant). The confidence margin provides a partial mechanistic account of \emph{why} certain examples are affected, but the \emph{existence} of per-example multiplicity is established independently.

The Jaccard similarity measure applied to the binary sets provides an interpretable metric of cross-architecture overlap that would be less transparent with continuous alternatives.

\location{\S4.3 (confidence margin paragraph), Appendix (Confidence Margin Analysis, Cross-Benchmark Extension, BoolQ Supplementary Analysis).}

\subsection*{W-ZOO (R1-W6): Reliance on a single model zoo}

\reviewer{All results depend on one model zoo; the controlled design limits external validity.}

\response{We acknowledge this trade-off explicitly. The controlled design of the zoo---identical training data, tokenizer, and hyperparameters across all configurations---is what allows us to isolate architectural effects from training confounds. This is stated in \S3.1 and in the Limitations section (``Architecture coverage'' paragraph), which now also notes: ``Each configuration in the model zoo corresponds to a single training run; we cannot disentangle architectural effects from seed-specific effects.'' The Conclusion calls for extending the analysis ``beyond these three architectures.'' We believe the controlled design is the appropriate choice for a first systematic study, while recognizing that confirmation across additional model families is needed. We additionally confirm that per-example multiplicity is comparable at 340M scale ($\alpha = 0.31$, $\bar{J} \approx 0.09$; Appendix~\S F), providing cross-scale consistency evidence within the zoo.}

\location{\S3.1, Limitations (``Architecture coverage''), \S5, Appendix~\S F.}

\subsection*{Q1: p-value rounding discrepancy (0.100 vs.\ 0.102)}

\reviewer{Why does the Bonferroni-adjusted p-value appear as both 0.100 and 0.102 in different places?}

\response{We have added a footnote to Table~1 explaining the discrepancy: ``Bonferroni-adjusted $p = 6 \times 2/120 = 0.100$ (exact); 0.102 arises from rounding raw $p$ to .017 before multiplying.'' The value 0.100 is the exact result; 0.102 is a rounding artifact. We now use 0.100 consistently throughout.}

\location{Table~1 footnote.}

\subsection*{Q5: Single checkpoint per configuration}

\reviewer{Each model zoo configuration is a single training run. Can you disentangle architecture from seed?}

\response{We cannot, and we now state this explicitly in the Limitations section: ``Each configuration in the model zoo corresponds to a single training run; we cannot disentangle architectural effects from seed-specific effects.'' This is a limitation of the zoo design. Disentangling the two would require multiple seeds per architecture-attention configuration, which is beyond the scope of the available zoo. We note, however, that the zoo's controlled design (shared data, hyperparameters) reduces other confounds, making architecture the most salient axis of variation. We additionally confirm that per-example multiplicity is comparable at 340M scale ($\alpha = 0.31$, $\bar{J} \approx 0.09$; Appendix~\S F), providing cross-scale consistency evidence within the zoo.}

\location{Limitations (``Architecture coverage'' paragraph), Appendix~\S F.}

\subsection*{W-REFS (R2-W2): Additional related work citations}

\reviewer{Several relevant references are missing from the related work discussion.}

\response{We have substantially expanded the citation coverage. In addition to Sellam et al.\ (ICLR 2022; MultiBERTs) and Fehlauer et al.\ (EMNLP 2025) added in Round~2, we now cite five additional references:}
\begin{itemize}[leftmargin=2em]
\item \textbf{Predictive Multiplicity paragraph}: Du et al.\ (2025; reconciliation under multiplicity), Haghighat et al.\ (2026; Rashomon-set resolution), and Gorbett \& Jana (2026; cross-model disagreement as label-free correctness signal).
\item \textbf{Hybrid Architectures paragraph}: Ali et al.\ (2025; implicit attention in Mamba) and Lahoti et al.\ (2026; Mamba-3).
\item \textbf{Appendix~D}: Ali et al.\ (2025) cited as methodological background for attention weight analysis.
\item \textbf{\S4.3}: Gorbett \& Jana (2026) cited for the practical relevance of per-example disagreement.
\end{itemize}

\location{\S2 (``Predictive Multiplicity'' and ``Hybrid Recurrent Architectures'' paragraphs), \S4.3, Appendix~D, \texttt{references.bib}.}

\end{document}
